% !TEX root = ../bb_AiRa_Kappaleet.tex

% ö = \%c3\%b6
% ä = \%C3\%A4

\xdef\sectiontitle{\IVsectiontitle} %aiheuttama
\TOCrefs

%%%%%%%%%%%%%%%
\begin{frame}[label=4_3_a]%[label=4_3_TOC1]
\frametitle{\texorpdfstring{\culine[\sectiontitleulinecolor]{\sectiontitle}}{\sectiontitle} }
\label{\TOCname}

\vspace{-3mm}

%\vspace{\chapterTOCvksip}%
\begin{itemize}[leftmargin=0.5cm,itemsep=\chapterTOCitemsep]
    \item[4.1]<1-> \hyperlink{4_1_a}{\IVaSubsectiontitle}
    \item[4.2]<1-> \hyperlink{4_2_a}{\IVbSubsectiontitle}	
    \item[4.3]<1-> \hyperlink{4_3_a}{\CDAlert[blue]{\IVcSubsectiontitle}}	
    \item[4.4]<1-> \hyperlink{4_4_a}{\IVdSubsectiontitle}
    \item[4.5]<1-> \hyperlink{4_5_a}{\IVeSubsectiontitle}
    \item[4.6]<1-> \hyperlink{4_6_a}{\IVfSubsectiontitle}
    \item[4.7]<1-> \hyperlink{4_7_a}{\IVgSubsectiontitle}
\end{itemize}

%\endofslide 
\end{frame}
%%%%%%%%%%%%%%%

\xdef\sectiontitle{\IVcSubsectiontitle} 


%%%%%%%%%%%%%%%
\begin{frame}
\frametitle{\texorpdfstring{\culine[\sectiontitleulinecolor]{\sectiontitle}}{\sectiontitle} }
\framesubtitle{Perusvuorovaikutukset}%On fundamental forces
%
\appear{}

\semismall
\begin{itemize}
	\item Ymmärryksemme perusvuorovaikutuksista on muuttunut paljon 1900-luvulta:%
	\appear{\xdef\loc{south}
\begin{tikzpicture}[remember picture,overlay] % anchor=south
    \node (A) [yshift=-0.35*\figYshift,xshift=0*\figXshift, anchor=\loc] at (current page.\loc) {\includegraphics[width=11cm]{Figures_booklet/SOM_11_Fundamental_particles_figures/SOM_Fundamental_Table_1_cropped.png}}; % 8cm max height
\end{tikzpicture}\CR[ref=h]{}%
}
\item<.->[] \begin{itemize}[leftmargin=0.2cm]
	\item Sähköoppi ja magnetismi \appear{yhdistettiin sähkömagnetismiksi} \appear{(Maxwell, 1873)}
	\item Sähkömagnetismi yhdistettiin kvanttimekaniikkaan ja suhteellisuusteoriaan (\CDAlert[blue]{kvanttielektrodynamiikka (QED)}) 1940–1950-luvuilla 
\item Sähkömagnetismi ja heikko vuorovaikutus yhdistettiin \CDAlert[blue]{sähköheikoksi voimaksi} 1960-luvulla. \appear{Näitä vuorovaikutuksia tarkastellaan usein erillisinä voimina}\appear{, koska "heikko voima" ~näyttäytyy vain tietyissä ytimien sisäisissä prosesseissa}
\item \CDAlert[red]{Vahvalle vuorovaikutukselle} muodostettiin kvanttikenttäteoria  \CDAlert[tunired]{kvanttikromodynamiikka (QCD)}
\item \CDAlert[fuchsia]{Gravitaatio} on säilynyt erillisenä voimanaan/\CDAlert[green]{vuorovaikutuksena}
\end{itemize}
	
\end{itemize}


\endofslide 
%\endofsection
\end{frame}
%%%%%%%%%%%%%%%

%%%%%%%%%%%%%%%
\begin{frame}
\frametitle{\texorpdfstring{\culine[\sectiontitleulinecolor]{\sectiontitle}}{\sectiontitle} }
\framesubtitle{Alkeishiukkaset}%On fundamental particles
\appear{}

\semismall
\begin{itemize}[itemsep=1.9mm]
	\item \appear{\xdef\loc{south east}%
\begin{tikzpicture}[remember picture,overlay] % anchor=south
    \node (A) [yshift=-0.25*\figYshift,xshift=1*\figXshift, anchor=\loc] at (current page.\loc) {\includegraphics[width=14.5cm]{Figures_booklet/SOM_11_Fundamental_particles_figures/Elementary_particles}};% 8cm max height
\end{tikzpicture}}%
\appear{\CDAlert[blue]{Fermionit kokevat voimia}:} \appear{\Alert[blue]{kvarkit}} \appear{ja \Alert[tuniblue]{leptonit}}
	\item \CDAlert[red]{Bosonit välittävät voimia}: \appear{gluonit}\appear{, W-} \appear{ja Z-bosonit}\appear{, fotonit}\appear{, gravitonit?}
	
	\item Jokaisella voiman kokevalla fermionilla\appear{, on voimaan liittyvä luontainen ominaisuus:}
\begin{itemize}
	\item \CDAlert[white!30!fuchsia]{Massa/energia}  <=> painovoima
	\item \CDAlert[fuchsia]{Varaus}\appear{, tai heikko varaus}\appear{ <=> sähköheikko vuorovaikutus}
	\item \CDAlert[black!30!fuchsia]{Värivaraus} <=>  vahva vuorovaikutus 
\end{itemize}
\end{itemize}

\endofslide 
%\endofsection
\end{frame}
%%%%%%%%%%%%%%%

%%%%%%%%%%%%%%%
\begin{frame}
\frametitle{\texorpdfstring{\culine[\sectiontitleulinecolor]{\sectiontitle}}{\sectiontitle} }
\framesubtitle{Alkeishiukkaset}%On fundamental particles
\appear{}

\begin{itemize}[itemsep=4mm]
	\item Alkeishiukkasten havainnointi tapahtuu yleensä epäsuorasti \appear{hyödyntämällä \CDAlert[red]{säilymislakeja} ja tunnettuja prosesseja} 
	\item Oletus joidenkin hiukkasten olemassaolosta perustuu \CDAlert[fuchsia]{standardimalliin} 
	\item Esimerkiksi kvarkkien\appear{ ja niitä sitovien gluonien välisistä vuorovaikutuksista}\appear{ ei ole vahvoja kokeellisia todisteita}
	
	\item Standardimallin ennustaman \CDAlert[blue]{Higgsin bosonin $\mathrm{H}^{0}$} löytyminen, ja sen massan määrittäminen 4.7.2012 ($125 \mathrm{\;GeV} / \mathrm{c}^{2}$)\appear{, on ollut yksi viimeaikaisista suurista löydöistä jotka ovat tarkentaneet kuvaamme aineen rakenteesta}
\implies{ Lupaus \CDAlert[fuchsia]{suuresta yhtenäisteoriasta} }
\end{itemize}

\endofslide 
%\endofsection
\end{frame}
%%%%%%%%%%%%%%%

%%%%%%%%%%%%%%%
\begin{frame}
\frametitle{\texorpdfstring{\culine[\sectiontitleulinecolor]{\sectiontitle}}{\sectiontitle} }
\framesubtitle{Vahva vuorovaikutus}%The strong force
%
\appear{}

\semismall
% 6.5 cm = midway, 10 cm = 3/4 
%\vspace{2.5mm}
\begin{minipage}{8.5cm}
\begin{itemize}[itemsep=1.7mm]
%	\item Nukleonien sisällä olevat kvarkit \appear{(protonit ja neutronit)} \appear{sitoutuvat toisiinsa vahvan vuorovaikutuksen kautta}%
\item Nukleonien sisällä olevat kvarkit sitoutuvat toisiinsa vahvan vuorovaikutuksen kautta protoneiksi ja neutroneiksi%
	\appear{\xdef\loc{north east}%
\begin{tikzpicture}[remember picture,overlay] % anchor=south
    \node (A) [yshift=0.75*\figYshift,xshift=\figXshift, anchor=\loc] at (current page.\loc) {\includegraphics[width=5cm]{Figures_booklet/SOM_11_Fundamental_particles_figures/SOM_Fundamental_particles_Figure_3.png}};% 8cm max height
\end{tikzpicture}\CR[ref=h]{}}%
	\item Kvarkkien ajatellaan olevan \CDAlert[blue]{fermionisia alkeishiuk-} \CDAlert[blue]{kasia jotka vuorovaikuttavat vahvan voiman} \CDAlert[blue]{kautta}\appear{, vaihtamalla gluoneita keskenään}
	
	\item \CDAlert[red]{Kuudenlaisia kvarkkeja} on olemassa \appear{(taulukko)}%up (u), down (d), strange (s), charm (c), bottom (b) and top (t)
\appear{\xdef\loc{south east}%
\begin{tikzpicture}[remember picture,overlay] % anchor=south
    \node (A) [yshift=-0.35*\figYshift,xshift=1*\figXshift, anchor=\loc] at (current page.\loc) {\includegraphics[width=8cm]{Figures_booklet/SOM_11_Fundamental_particles_figures/SOM_Fundamental_Table_1_cropped2.png}};% 8cm max height
\end{tikzpicture}}%\CR[]{}
\item Jokaisen spin on $\frac{1}{2}$\appear{, ja varaus joko} \appear{$-\Frac{1}{3} e$} \appear{tai $+\Frac{2}{3} e$}

\end{itemize}
\end{minipage}

% 6.5 cm = midway, 10 cm = 3/4 
\vspace{1.7mm}
\begin{minipage}{5.8cm}
\begin{itemize}[itemsep=1.7mm]
	\item \CDAlert[fuchsia]{Protoni} koostuu \CDAlert[fuchsia]{uud}-kvarkeista: 
	\[
	\text{varaus}  \equal  \plus \frac{2}{3} \plus \frac{2}{3} \plus \appear{(} \minus \frac{1}{3}\appear{)}\appear{=1}
	\]

\item \CDAlert[gray]{Neutroni} koostuu \CDAlert[gray]{udd}-kvarkeista:
\[
\text{varaus}  \equal  \plus \frac{2}{3} \plus \appear{(} \minus \frac{1}{3}\appear{)} \plus \appear{(} \minus \frac{1}{3}\appear{)} \equal \appear{0}
\]
\item Hiukkasia, jotka koostuvat kvarkeista kutsutaan \CDAlert[fuchsia]{hadroneiksi}
\end{itemize}
\end{minipage}
%\vspace{2.5mm}

\endofslide 
%\endofsection
\end{frame}
%%%%%%%%%%%%%%%

%%%%%%%%%%%%%%%
\begin{frame}
\frametitle{\texorpdfstring{\culine[\sectiontitleulinecolor]{\sectiontitle}}{\sectiontitle} }
\framesubtitle{Kvarkit/kvarkkien vankeus}%Quark/Color confinement
%Quark Confinement
\semismall
\appear{}

\begin{itemize}[itemsep=1.7mm]
	\item Yksittäisiä kvarkkeja ei ole havaittu. \appear{Ajatellaan että \Alert[tunired]{vahva voima voimistuu} etäisyyden kasvaessa} \appear{(toisin kuin sähköinen voima).} \appear{On lähestulkoon mahdotonta erottaa kvarkkeja toisistaan}\appear{, mitä ilmiötä kutsutaan} \appear{$\Leftrightarrow$ \CDAlert[red]{kvarkkien vankeudeksi} }
	
	\item Kun kvarkkeja yritetään erottaa niitä pommittamalla\appear{, systeemiin tuotu energia} \appear{vain \Alert[tunired]{synnyttää uusia kvarkki–antikvarkki-pareja}}
\appear{\xdef\loc{south east}%
\begin{tikzpicture}[remember picture,overlay] % anchor=south
    \node (A) [yshift=-0.35*\figYshift,xshift=\figXshift, anchor=\loc] at (current page.\loc) {\includegraphics[scale=0.75]{Figures_booklet/SOM_11_Fundamental_particles_figures/SOM_Fundamental_quark_confinement.pdf}};% 8cm max height
\end{tikzpicture}}

\end{itemize}

\vspace{1.7mm}
\begin{minipage}{8.2cm}
\begin{itemize}[itemsep=1.7mm]
	\item Kvarkeista on saatu tietoa epäelastisen sironnan avulla käyttäen suurenergisia elektroneita \appear{$E=20 \mathrm{\;GeV}$} \appear{\& $p=20 \mathrm{\;GeV} / \mathrm{c}$} \appear{(eli $\lambda \sim{10^{-16}} \mathrm{\;m}$)} \appear{(esim.\ SLAC, 1973)}
	\item Kokeiden perusteella \appear{kvarkit liikkuvat nukleonin sisällä suhteellisen vapaasti}
	\item Nukleonin muodostavan kolmen kvarkin yhteenlaskettu massa \appear{on huomattavasti alhaisempi kuin nukleonin massa.} \appear{\CDAlert[fuchsia]{Puuttuvan energian/massan} ajatellaan olevan varastoitunut nukleonin sisällä olevaan \CDAlert[fuchsia]{energeettiseen gluonipilveen}}
\end{itemize}
\end{minipage}

\endofslide 
%\endofsection
\end{frame}
%%%%%%%%%%%%%%%

%%%%%%%%%%%%%%%
\begin{frame}
\frametitle{\texorpdfstring{\culine[\sectiontitleulinecolor]{\sectiontitle}}{\sectiontitle} }
\framesubtitle{Gluonit}
\appear{}

\seminormal

\begin{itemize}[itemsep=1.6mm]
	\item \CDAlert[red]{Gluoni} on vahvan voiman \CDAlert[red]{välittäjähiukkanen}
	\item Kvarkit vaihtavat gluoneita jatkuvasti
\end{itemize}

% 6.5 cm = midway, 10 cm = 3/4 
\vspace{1.6mm}
\begin{minipage}{9.6cm}
\begin{itemize}[itemsep=1.6mm]
	\item Gluonien spin on 1, eikä niillä ole lepomassaa

\item Massattomalla välittäjähiukkasella\appear{ on ääretön kantama} \appear{((kantama} $\approx \frac{\hbar}{c} \frac{1}{m}$\appear{))} 

\item \CDAlert[blue]{Gluonit siirtävät} voimalle ominaista ominaisuutta\appear{, eli \CDAlert[blue]{värivarausta}}
\appear{(\CDAlert[tuniblue]{Toisin kuin fotonit} jotka ovat sähköisesti neutraaleita)}
\implies{ Koska gluoneilla on väri, niiden efektiivinen kantama \\
 		on erittäin lyhyt}

\item Vahvan voiman aiheuttamaa \CDAlert[red]{potentiaalia}\\ \CDAlert[red]{on approksimoitu}\appear{, ja kokeellisesti mitattu lyhyillä etäisyyksillä:}\vspace{-3.5mm}
\[
\qquad V_{QCD}(r)  \equal  \minus \frac{4 \alpha_{s}}{3 r}  \plus \appear{k r}
\]
\item[] missä $\alpha_{S} \appear{=\text{vahvan voiman kytkentävakioa}} \appear{\Approx0,3}$ \appear{and $k=1 \mathrm{\;GeV} / \mathrm{fm}$}
\end{itemize}

\end{minipage}

\appear{\xdef\loc{south east}
\begin{tikzpicture}[remember picture,overlay] % anchor=south
    \node (A) [yshift=-0.25*\figYshift,xshift=\figXshift, anchor=\loc] at (current page.\loc) {\includegraphics[scale=0.7]{Figures_booklet/SOM_11_Fundamental_particles_figures/Quark_potential.pdf}}; % 8cm max height
\end{tikzpicture}}

\endofslide 
%\endofsection
\end{frame}
%%%%%%%%%%%%%%%

%%%%%%%%%%%%%%%
\begin{frame}
\frametitle{\texorpdfstring{\culine[\sectiontitleulinecolor]{\sectiontitle}}{\sectiontitle} }
\framesubtitle{Värivaraus}
\appear{}

\seminormal
\begin{itemize}
	\item \CDAlert[fuchsia]{Värivaraus} \appear{(tai lyhyesti väri)} \appear{on vahvaan voimaan liittyvä ominaisuus} 
	\appear{\xdef\loc{south east}%
\begin{tikzpicture}[remember picture,overlay] % anchor=south
    \node (A) [yshift=-0.25*\figYshift,xshift=\figXshift, anchor=\loc] at (current page.\loc) {\includegraphics[scale=0.6]{Figures_booklet/SOM_11_Fundamental_particles_figures/Color_charges.pdf}};% 8cm max height
\end{tikzpicture}}%
	\vspace{2mm}
	\begin{itemize}[itemsep=2.2mm,leftmargin=1.5mm]
		\item \Alert[white!30!fuchsia]{Kvarkilla} on värivaraus: \appear{\textcolor{red}{\Alert[red]{punainen} (r)}}\appear{, \textcolor{green}{\CDAlert[green]{vihreä} $(g)$}}\appear{, ja \textcolor{blue}{\CDAlert[blue]{sininen} $(b)$}}
	\item \Alert[black!30!fuchsia]{Antikvarkeilla} on antiväri: \appear{\textcolor{antired}{\CDAlert[antired]{antipunainen} $(\bar{r})$}}\appear{, \textcolor{antigreen}{\CDAlert[antigreen]{antivihreä} $(\bar{g})$}}\appear{, ja \textcolor{antiblue}{\CDAlert[antiblue]{antisininen} $(\bar{b})$}}
	
	\item Värit siis kuvaavat vain hiukkasten kvanttimekaanista ominaisuutta \appear{(eivät ole näkyviä värejä\,!!)}
	\end{itemize}
\end{itemize}

% 6.5 cm = midway, 10 cm = 3/4 
%\vspace{2.5mm}
\begin{minipage}{5.1cm}
\begin{itemize}	
\item[] \begin{itemize}[itemsep=2.2mm,leftmargin=1.5mm]
	\item Kuten sähköinen varaus\appear{, myös \CDAlert[fuchsia]{väri säilyy}}
\item Gluonit itsessään kuljettavat väriä ja antiväriä \appear{((on 8 erilaista gluonia, SU(3)\,))}
\item \CDAlert[red]{Gluonit vuorovaikuttavat} \CDAlert[red]{myös toisten gluonien kanssa!}
\end{itemize}
\end{itemize}
\end{minipage}
%\vspace{2.5mm}

\endofslide 
%\endofsection
\end{frame}
%%%%%%%%%%%%%%%

%%%%%%%%%%%%%%%
\begin{frame}
\frametitle{\texorpdfstring{\culine[\sectiontitleulinecolor]{\sectiontitle}}{\sectiontitle} }
\framesubtitle{Miksi väri tarvittiin/otettiin käyttöön?}%Why color was needed/introduced?
\appear{}

\seminormal
%Why was color introduced (needed) in the first place?s
\begin{itemize}
	\item Hadronit koostuvat \CDAlert[blue]{fermionista kvarkeista} \appear{ja ovat värineutraaleita}
	
	\item \CDAlert[tuniblue]{Kieltosääntö} estää kahta fermionia olemasta samassa tilassa \appear{(joilla olisi täysin saman kvanttiluvut)} 
	\implies{ Hadronien sisällä olevat kvarkit eivät voi olla samassa tilassa } 
	\implies{ Värin ominaisuus (3 erilaista arvoa) esiteltiin vuonna 1964 O. W. Greenbergin toimesta} 
	\item Ilman väriä\appear{, esim.\ hadronit} $\appear{\Delta^{++}=uuu} \appear{~\&~ \Delta^{-}=d d d}$ \appear{olisivat samassa kvanttitilassa}

\item Esim.\ nukleonin sisäisten kolmen kvarkin värit poikkeavat toisistaan

\item \Alert[red]{Mesonit} koostuvat \CDAlert[tunired]{kvarkki–antikvarkki-parista} 
\implies{Värineutraalisuus vaatii että \CDAlert[green]{väri} \appear{ja \CDAlert[antigreen]{antiväri} muodostavat parin}}

\item \CDAlert[fuchsia]{Kokeellisesti havaitut} energeettisten hiukkasten törmäyksistä sinkoutuvat suihkut \CDAlert[fuchsia]{implikoivat gluonien ja värin olevan olemassa} 
%\item It can be deducted from the collision products, that sometimes individual gluons have been produced immediately after collision
\end{itemize}

\endofslide 
%\endofsection
\end{frame}
%%%%%%%%%%%%%%%

%%%%%%%%%%%%%%%
\begin{frame}
\frametitle{\texorpdfstring{\culine[\sectiontitleulinecolor]{\sectiontitle}}{\sectiontitle} }
\framesubtitle{Värivaraus ja kvarkkien vankeus}%Color charge and color confinement
\appear{}

\semismall
%Electrical charge vs. Color charge
\begin{itemize}
%	\item How can we interpret the difference that for electroweak force the property electrical charge is free and 'abundant' in the particles, but for the specific property of strong force, color, all the hadrons are color-neutral? 
	
	\item Sähköinen varaus ja värivaraus käyttäytyvät eri tavalla: \vspace{2mm}
	
	\begin{itemize}
	 \item Sähköisesti varautuneita hiukkasia löytyy paljonkin yksittäisinä hiukkasina \appear{(esim.\ elektroni)}
	 \item Hadronit ovat värineutraaleita\appear{, värillistä hiukkasta ei ole onnistuttu synnyttämään!}
	 \item \CDAlert[fuchsia]{Miksei vapaita värivarauksia löydy?}
\end{itemize}

	\item Yksinkertaistettu selitys on \appear{että voimien suuruusluokat eroavat toisistaan huomattavasti:} \appear{Vahvasti vuorovaikuttavat hiukkaset vetävät toisiaan voimakkaasti puoleensa}
	
\end{itemize}

% 6.5 cm = midway, 10 cm = 3/4 
\vspace{2.5mm}
\begin{minipage}{8.2cm}
\begin{itemize}
	\item Mutta vuorovaikutusten voimakkuudet eivät selitä miksi \CDAlert[red]{yksittäistä kvarkkia ei} \CDAlert[red]{ole kokeellisesti havaittu!}
	
	\item Vastaus piilee \CDAlert[tuniblue]{gluonien vuorovaikutuksissa} \\\CDAlert[tuniblue]{toisten gluonien kanssa}\appear{, mikä synnyttää \CDAlert[blue]{voiman} \CDAlert[blue]{mikä kasvaa}} \appear{värillisten hiukkasten etäisyyden kasvaessa}
\appear{\xdef\loc{south east}%
\begin{tikzpicture}[remember picture,overlay] % anchor=south
    \node (A) [yshift=-0.25*\figYshift,xshift=\figXshift, anchor=\loc] at (current page.\loc) {\includegraphics[scale=0.75]{Figures_booklet/SOM_11_Fundamental_particles_figures/SOM_Fundamental_quark_confinement.pdf}};% 8cm max height
\end{tikzpicture}}
\end{itemize}

\end{minipage}
%\vspace{2.5mm}

\endofslide 
%\endofsection
\end{frame}
%%%%%%%%%%%%%%%

%%%%%%%%%%%%%%%
\begin{frame}
\frametitle{\texorpdfstring{\culine[\sectiontitleulinecolor]{\sectiontitle}}{\sectiontitle} }
\framesubtitle{Hadronien luokittelu}%Categories of hadrons
\appear{}

\seminormal
%Categories of Hadrons
\begin{itemize}[itemsep=1.9mm]
	\item Koska on olemassa kuusi kvarkkia\appear{/antikvarkkia}%
	\appear{\xdef\loc{south east}%
\begin{tikzpicture}[remember picture,overlay] % anchor=south
    \node (A) [yshift=-0.35*\figYshift,xshift=\figXshift, anchor=\loc] at (current page.\loc) {\includegraphics[width=8cm]{Figures_booklet/SOM_11_Fundamental_particles_figures/SOM_Fundamental_Table_1_cropped2.png}};% 8cm max height
\end{tikzpicture}\CR[ref=h]{}}%
\appear{, monia hadroneita voidaan muodostaa}

	\item Taulukko 12.2. (seuraava kalvo) listaa yleisimmät hadronit
	
	\item \CDAlert[fuchsia]{Protoni on ainoa vakaa hadroni} \appear{vapaan neutronin elinaika on $889 \mathrm{\;s}$)}
	%free neutron is $889 \mathrm{sec})$. Note, that instead of lifetime, an energy range is given assuming $\Delta E \approx \hbar / \tau$, where $\tau$ is the time interval where the particle has been observed to exist
	\item \CDAlert[blue]{Baryonit} koostuvat  \CDAlert[blue]{kolmesta kvarkista}: \appear{esim.\ Protoni (uud) ja neutroni (udd)}
	
	\item \CDAlert[red]{Mesonit} koostuvat \CDAlert[red]{kvarkista ja antikvarkista}: \appear{$\pi^{+}$-mesoni} \appear{($u\bar{d}$)}\appear{, ja $\pi^{-}$-mesoni ($\bar{u}d$)}

\end{itemize}

% 6.5 cm = midway, 10 cm = 3/4 
\vspace{2.5mm}
\begin{minipage}{5.5cm}
\begin{itemize}
	\item Elinajan sijasta\appear{ usein annetaan \CDAlert[fuchsia]{energiaväli} olettaen} $\appear{\Delta E} \approx \appear{\hbar / \tau}$\appear{, missä $\tau$ on aikaväli jona hiukkasen oletetaan olleen olemassa}
\end{itemize}
\end{minipage}

\endofslide 
%\endofsection
\end{frame}
%%%%%%%%%%%%%%%

%%%%%%%%%%%%%%%
\begin{frame}
\frametitle{\texorpdfstring{\culine[\sectiontitleulinecolor]{\sectiontitle}}{\sectiontitle} }
\framesubtitle{Yleisimmät hadronit}
\appear{}

\appear{\xdef\loc{south}
\begin{tikzpicture}[remember picture,overlay] % anchor=south
    \node (A) [yshift=-0.25*\figYshift,xshift=0*\figXshift, anchor=\loc] at (current page.\loc) {\includegraphics[width=15cm]{Figures_booklet/SOM_11_Fundamental_particles_figures/SOM_Fundamental_Table_2.png}}; % 8cm max height
\end{tikzpicture}\CR[ref=h]{}}

\endofslide 
%\endofsection
\end{frame}
%%%%%%%%%%%%%%%

%%%%%%%%%%%%%%%
\begin{frame}
\frametitle{\texorpdfstring{\culine[\sectiontitleulinecolor]{\sectiontitle}}{\sectiontitle} }
\framesubtitle{Luontaiset ominaisuudet}%Intrinsic properties
\appear{}

\seminormal
\begin{itemize}
	\item Kaikki kvarkit ovat spinin $\Frac{1}{2}$ fermioneita\appear{, tällöin kulmaliikemäärän yhteenlaskusääntöjen nojalla}\appear{, kaikki baryonit} \appear{(kolme kvarkkia)} \appear{ovat joko spinin $\Frac{1}{2}$} \appear{tai $\Frac{3}{2}$ hiukkasia}

\item \CDAlert[red]{Isospin} \vspace{2mm}

\begin{itemize}
	\item Taulukko 12.2 näyttää\appear{ että joillakin hadroneilla näyttää olevat samanlainen kvarkkikoostumus} \appear{ja sama spin} 
	\item Toteuttaakseen kieltosäännön\appear{, uusi kvanttiluku täytyy ottaa käyttöön}
	\item Tätä kvanttilukua kutsutaan isospiniksi $I$\appear{, ja se on vielä abstraktimpi kuin tavallinen spin}\appear{, mutta sen \CDAlert[tunired]{kolmas komponentti $I_{3}$} kvantisoituu samalla tavalla kuin spin } 
	
	\item Ylös- ja alas-kvarkin isospinit ovat $I=\frac{1}{2}$\appear{, oudon kvarkin $I=0$}
\end{itemize}
\item \CDAlert[blue]{Outous} \vspace{2mm}
	\begin{itemize} 
	\item Oudon kvarkin outous on $s=-1$. \appear{Eli esim.\  hadroni $\Omega^{-}(sss): \quad s=-3$}
	\end{itemize}
\end{itemize}

\endofslide 
%\endofsection
\end{frame}
%%%%%%%%%%%%%%%

%%%%%%%%%%%%%%%
\begin{frame}
\frametitle{\texorpdfstring{\culine[\sectiontitleulinecolor]{\sectiontitle}}{\sectiontitle} }
\framesubtitle{Vahvan voiman jäännösvoima (ydinvoima)}%The residual strong force
\appear{}

\begin{itemize}[itemsep=2.2mm]
	\item Sähköisesti neutraalit molekyylit voivat polarisoitua sähköisesti\appear{, ja alkaa vuorovaikuttamaan sähköisesti keskenään} 
	\item Vesimolekyylit ovat hyvä esimerkki \appear{muodostaessaan jäätä} 
	\item Vastaavasti värineutraalit hiukkaset voivat vuorovaikuttaa keskenään vahvan voiman kautta \implies{ \CDAlert[blue]{Tämä vuorovaikutus/vetovoima sitoo nukleonit toisiinsa atomien ytimissä!} }
	
\item Tämä niin sanottu \CDAlert[blue]{vahva jäännösvoima} ei oikeasti ole oma voimansa\appear{, vaan \CDAlert[black!30!blue]{oikean vahvan} \CDAlert[tuniblue]{vuovovaikutuksen} \CDAlert[tuniblue]{ilmentymä} }

\item Protoni ja neutroni \appear{(värineutraaleja)} \appear{vaihtavat pioneja} \appear{(kahden kvarkin mesoneita) }
\item \CDAlert[red]{Mesonin massa} selittää jäännösvoiman \CDAlert[red]{lyhyen kantaman} 
\item $1 \mathrm{\;fm}$ kantama antaa välittäjähiukkasen massaksi $197 \mathrm{\;Mev} / \mathrm{c}^{2}$\appear{, mitattu arvo on $140 \mathrm{\;Mev} / \mathrm{c}^{2}$}
\end{itemize}

\endofslide 
%\endofsection
\end{frame}
%%%%%%%%%%%%%%%

%%%%%%%%%%%%%%%
\begin{frame}
\frametitle{\texorpdfstring{\culine[\sectiontitleulinecolor]{\sectiontitle}}{\sectiontitle} }
\framesubtitle{Sähköheikko vuorovaikutus}%Electroweak force
\appear{}

\seminormal
\begin{itemize}[itemsep=1.8mm]
	\item Hiukkaset joilla on \CDAlert[red]{varaus}/\CDAlert[blue]{heikko varaus} \appear{(joskus kutsutaan myös '\CDAlert[blue]{makuvaraukseksi}')}\appear{, voivat vuorovaikuttaa \Alert[red]{sähkö}\Alert[blue]{heikon} voiman kautta }
	\implies{ Sekä kvarkit että leptonit kokevat sähköheikkoja voimia}
	
	 \item \CDAlert[fuchsia]{Fermioniset leptonit} ovat alkeishiukkasia \CDAlert[fuchsia]{joilla ei ole värivarausta}\appear{, eivätkä ne siis vuorovaikuta vahvan voiman kautta}

\item Leptoneita on olemassa \CDAlert[fuchsia]{kolme sukupolvea}: \appear{elektronit}\appear{, myonit} \appear{ja tauonit} 
\item Jokaisella sukupolvella on oma leptoninsa\appear{, antileptoni}\appear{, neutriino} \appear{ja antineutriino}

\item Elektroni, myoni ja tauoni \appear{käyttäytyvät hyvinkin samanlaisesti}\appear{, niillä on vain erisuuret massat}

\item \CDAlert[blue]{Sähköheikko vuorovaikutus hallitsee} \appear{etäisyyksillä jotka ovat ydintä suurempia} \appear{mutta kosmologisia etäisyyksiä pienempiä} 
\item Sen ansiosta aineen eri faasit muodostuvat\appear{, sekä \CDAlert[blue]{fyysiset ja kemialliset}}
\end{itemize}

\endofslide 
%\endofsection
\end{frame}
%%%%%%%%%%%%%%%

%%%%%%%%%%%%%%%
\begin{frame}
\frametitle{\texorpdfstring{\culine[\sectiontitleulinecolor]{\sectiontitle}}{\sectiontitle} }
\framesubtitle{Sähköheikon voiman välittäjähiukkaset}%Mediating particles of electroweak force
\appear{}

%Mediating Particles of Electroweak Force
\begin{itemize}
	\item Sähköheikon voiman välittäjähiukkanen on massaton fotoni \appear{ja kolme massallista bosonia}\appear{, $\mathrm{W}^{+}$}\appear{, $\mathrm{W}^{-}$} \appear{ja $\mathrm{Z}^{0}$ }
	
	\item Sähkömagneettista ja heikkoa vuorovaikutusta pidetään siltikin erillisinä voimina\appear{, joskin 1970-luvulla}\appear{ molemmat yhdistettiin kvanttielektrodynamiikan (QED) alle}

\item \CDAlert[fuchsia]{Välittäjähiukkasten massoissa oleva epäsymmetria} selitettiin, taas, analogian avulla 
\item Hyrrä voi pyöriä symmetrisesti (ei prekessoi) suurilla energioilla\appear{, mutta kaatua spotaanisti matalammilla energioilla synnyttäen epäsymmetrisen tilanteen }
\item \CDAlert[fuchsia]{Symmetria} olisi näkyvissä vain erittäin suurilla energioilla $1 \mathrm{\;TeV}$\appear{, jota ei ole vielä havaittu}
\end{itemize}

\endofslide 
%\endofsection
\end{frame}
%%%%%%%%%%%%%%%

%%%%%%%%%%%%%%%
\begin{frame}
\frametitle{\texorpdfstring{\culine[\sectiontitleulinecolor]{\sectiontitle}}{\sectiontitle} }
\framesubtitle{Gravitaatiovoima}
\appear{}

\seminormal
\begin{itemize}[itemsep=2mm]
	\item Kaikki hiukkaset kokevat gravitaatiovoiman\appear{, koska sen lähde on \CDAlert[blue]{"gravitaatiovaraus" ~eli massa}} \appear{((... \Alert[fuchsia]{yleisessä suhteellisuusteoriassa stressi–energia-tensori}...))}
	
	\item \appear{\CDAlert[blue]{Painovoima on}} \appear{${\sim}10^{-38}$ kertaa heikompi kuin vahva voima} 
	\implies{Painovoima ei vaikuta alkeishiukkasten, atomien/molekyylien muodostumiseen}

\item Vuorovaikuttajahiukkanen olisi \CDAlert[tuniblue]{vielä löytymätön} \CDAlert[blue]{gravitoni}

\item Sen oletetaan olevan sähköisesti neutraali\appear{, \CDAlert[blue]{massaton}} \appear{ja omaavan \CDAlert[blue]{spinin 2}}

\item Menetelmiä ja laitteita on kehitteillä gravitonin löytämiseksi

\item \CDAlert[red]{Gravitaatioaallot} ovat puolestaan häiriöitä aika-avaruuden rakenteessa ("kaarevuudessa") joita kiihtyvät massat aiheuttavat\appear{, jotka etenevät valonnopeudella aaltoina lähteestä poispäin} 

\item \CDAlert[red]{2017}, Weiss, Thorne \& Barich saivat \CDAlert[red]{Nobelin fysiikanpalkinnon} rooleistaan \CDAlert[tunired]{gravitaatioaaltojen havainnoinnissa}

\end{itemize}

\endofslide 
%\endofsection
\end{frame}
%%%%%%%%%%%%%%%

%%%%%%%%%%%%%%%
\begin{frame}
\frametitle{\texorpdfstring{\culine[\sectiontitleulinecolor]{\sectiontitle}}{\sectiontitle} }
\framesubtitle{Standardimalli}
\appear{}

\seminormal 
\begin{itemize}
	\item \CDAlert[fuchsia]{Standardimalli} on tämänhetkinen hyväksytty teoria kuvaamaan alkeishiukkasia:
\end{itemize}

% 6.5 cm = midway, 10 cm = 3/4 
\vspace{1.5mm}
\begin{minipage}{8.6cm}
\begin{itemize}
\item<.->[] \begin{itemize}
		\item Hiukkasten rakenteen \CDAlert[red]{kvarkkimalli}
		\item Sähkömagneettisen ja heikon voiman yhdistäminen \appear{$\oldrightarrow$ sähköheikko teoria}
		\item \CDAlert[blue]{sähköheikon ja vahvan voiman yhdistäminen} \appear{(kvanttikromodynamiikka (QCD))} \appear{$\oldrightarrow$ suuri yhtenäisteoria (GUT)}
		\item Selittää alkeishiukkasten ominaisuudet
		\item Selittää vuorovaikutusten luonteen
	\end{itemize}
	\vspace{0mm}
	\item Joitain \CDAlert[fuchsia]{avoimia kysymyksiä}: \vspace{1.5mm}
\begin{itemize}\semismall
	\item \CDAlert[red]{Neutriinon lepomassa?}
\item \CDAlert[green]{Higgsin bosonin} ominaisuudet?
\item \CDAlert[blue]{Gravitonia ei ole havaittu}
\item Galaksien pyörimiskäyrät\appear{, \Alert[fuchsia]{pimeä aine}/\Alert[red]{energia}?}
\end{itemize}

\end{itemize}
\end{minipage}
%\vspace{2.5mm}

\appear{\xdef\loc{south east}
\begin{tikzpicture}[remember picture,overlay] % anchor=south
    \node (A) [yshift=-0.25*\figYshift,xshift=0*\figXshift, anchor=\loc] at (current page.\loc) {\includegraphics[scale=0.5]{Figures_booklet/SOM_11_Fundamental_particles_figures/Separation_of_forces_schematic.pdf}}; % 8cm max height
\end{tikzpicture}}


%\endofslide 
\endofsection
\end{frame}
%%%%%%%%%%%%%%%

